\subsection{Characters and conjugacy classes}

Let \(\varrho\) be a unitary representation of G in a finite-dimensional Hilbert space E. The character of \(\varrho\) satisfies \(|\chi_{\varrho}| \leqslant \dim(E)\) . Let \((e_i)_{i \in I}\) be an orthonormal basis of E; the character \(\chi_{\varrho}\) is the sum of the diagonal matrix coefficients \(g \mapsto \langle e_i | \varrho(g)e_i \rangle\) of \(\varrho\) . In particular, the character of \(\varrho\) is a G-finite function, and if \(\varrho\) is irreducible, then \(\chi_{\varrho} \in \varrho_{\mathrm{G}}(\varrho)\) .

We have the following formulas

\[
\chi_ {\varrho_ {1} \oplus \varrho_ {2}} = \chi_ {\varrho_ {1}} + \chi_ {\varrho_ {2}}, \quad \chi_ {\varrho_ {1} \otimes \varrho_ {2}} = \chi_ {\varrho_ {1}} \chi_ {\varrho_ {2}}, \quad \chi_ {\breve {\varrho}} = \chi_ {\overline {{\varrho}}} = \overline {{\chi}} _ {\varrho}
\]

(cf. A, VIII, p. 388--389).

PROPOSITION 8. --- We have

\begin{enumerate}
\def\labelenumi{(\arabic{enumi})}
\tightlist
\item
  \(\chi_{\pi}*\chi_{\sigma} = 0\) for all \(\pi, \sigma\) belonging to \(\widehat{\mathrm{G}}\) , \(\pi \neq \sigma\) ,
\end{enumerate}

\[
\chi_ {\pi} * \chi_ {\pi} = \frac {1}{\dim (\pi)} \chi_ {\pi} \quad\text{for every }\pi\in\widehat {\mathrm{G}}.\tag{2}
\]

Let \(\pi\) and \(\sigma\) be irreducible representations of G. We have

\[
\dim (\pi) \left(\chi_ {\pi} * \chi_ {\sigma}\right) = \dim (\pi) \boldsymbol {\gamma} _ {G} \left(\chi_ {\pi}\right) \chi_ {\sigma}
\]

(lemma 4 of V, p. 407). The function \(\dim(\pi)\gamma_{\mathrm{G}}(\chi_{\pi})\chi_{\sigma}\) is the orthogonal projection of \(\chi_{\sigma}\) onto the \(\overline{\pi}\)-isotypic component of \(\gamma_{\mathrm{G}}\) (cor. 2 of V, p. 465), that is, onto \(\varrho_{\mathrm{G}}(\pi)\) (cor. 4 of V, p. 463). Since \(\chi_{\sigma}\) belongs to \(\varrho_{\mathrm{G}}(\sigma)\) , the result follows from theorem 1 of V, p. 462.

Lemma 2. --- The graph of the equivalence relation « \(x \in G\) and \(y \in G\) and \(x\) is conjugate to \(y\) » in G is closed.

Indeed, this graph is the image of the continuous map from the compact space G × G into itself defined by \((x, y) \mapsto (x, yxy^{-1})\) .

We denote by G \(^{\sharp}\) the space of conjugacy classes of G equipped with the quotient topology; it is a compact space by lemma 2 and TG, I, p. 78, prop. 8. Let \(\varpi\) : G → G \(^{\sharp}\) be the canonical projection. The map from \(\mathcal{C}(G^{\sharp})\) into \(\mathcal{C}(G)\) defined by \(f \mapsto f \circ \varpi\) identifies the stellar algebra \(\mathcal{C}(G^{\sharp})\) with the stellar subalgebra of \(\mathcal{C}(G)\) formed by continuous central functions.

The measures on G \(^{\sharp}\) are identified with the central measures on G (V, p. 402, def. 1).

The linear form on \(\mathcal{C}(\mathrm{G}^{\sharp})\) defined by \(f\mapsto\int_{\mathrm{G}}f\) is then a positive measure of mass 1 on \(\mathrm{G}^{\sharp}\) , which is denoted \(\mu^{\sharp}\) . For every \(p\in[1,+\infty]\) , we denote by \(\mathcal{L}^{p}(\mathrm{G}^{\sharp})\) (resp. \(\mathrm{L}^{p}(\mathrm{G}^{\sharp})\) ) the space \(\mathcal{L}_{\mathbf{C}}^{p}(\mathrm{G}^{\sharp},\mu^{\sharp})\) (resp. \(\mathrm{L}_{\mathbf{C}}^{p}(\mathrm{G}^{\sharp},\mu^{\sharp})\) ). We identify \(\mathrm{L}^{p}(\mathrm{G}^{\sharp})\) with the closure of \(\mathcal{C}(\mathrm{G}^{\sharp})\) in \(\mathrm{L}^{p}(\mathrm{G})\) ; in particular, it is a closed subspace of \(\mathrm{L}^{p}(\mathrm{G})\) .

We also denote by \(\Theta(G^{\sharp}) = \Theta(G) \cap \mathcal{C}(G^{\sharp})\) the space of central matrix coefficients of G. It is a unital involutive subalgebra of \(\mathcal{C}(G^{\sharp})\) .

When G is a Lie group, the space \(\Theta(G^{\sharp})\) is also denoted \(Z\Theta(G)\) in LIE, IX, p. 71.

Lemma 3. --- Let \(\pi\) be an irreducible unitary representation of G. The vector space \(\mathrm{L}^2 (\mathrm{G}^\sharp)\cap \varrho_{\mathrm{G}}(\pi)\) is one-dimensional and generated by the character of \(\pi\) .

Consider the unitary representation \(\sigma\) of G in the space \(\varrho_{\mathrm{G}}(\pi)\) defined by \(\sigma(g) = \varrho_{\mathrm{G}}(g, g)\) . Since the unitary representation \(\varrho_{\mathrm{G}}(\pi)\) of G × G is isomorphic to \(\overline{\pi} \boxtimes \pi\) (prop. 5 of V, p. 462), its character is given by

\[
\chi_ {\sigma} (g) = \chi_ {\overline {{\pi}} \boxtimes \pi} (g, g) = | \chi_ {\pi} (g) | ^ {2}
\]

for all \(g \in G\) . The space \(\mathrm{L}^2 (\mathrm{G}^\sharp) \cap \varrho_{\mathrm{G}}(\pi)\) is the subspace of invariant elements of this representation, and its dimension is equal to

\[
\int_ {\mathrm{G}} \chi_ {\sigma} (g) d \mu (g) = \int_ {\mathrm{G}} | \chi_ {\pi} (g) | ^ {2} d \mu (g) = 1
\]

(cor. 3 of V, p. 466 and cor. 3 of V, p. 459). Since the character of \(\pi\) is a nonzero element of \(L^2(G^\sharp) \cap \varrho_G(\pi)\) , it is a basis of this space.

PROPOSITION 9. --- The family \((\chi_{\pi})_{\pi \in \widehat{\mathbf{G}}}\) of characters of irreducible unitary representations of G is an orthonormal basis of \(\mathrm{L}^2 (\mathrm{G}^\sharp)\) .

By the Peter--Weyl theorem (th. 1 of V, p. 462), the closed subspace \(L^{2}(G^{\sharp})\) of \(L^{2}(G)\) formed by the invariant elements of the unitary representation \(g \mapsto \varrho_{\mathrm{G}}(g, g)\) of G is the Hilbert sum of the subspaces \(L^{2}(G^{\sharp}) \cap \varrho_{\mathrm{G}}(\pi)\) for \(\pi \in \widehat{G}\) . The proposition therefore follows from the preceding lemma and cor. 3 of V, p. 459.

COROLLARY 1. --- The vector space \(\Theta(\mathrm{G}^{\sharp})\) is dense in \(\mathcal{C}(\mathrm{G}^{\sharp})\) and in \(\mathrm{L}^{2}(\mathrm{G}^{\sharp})\) .

The second assertion follows from prop. 9. As for the first assertion, consider the linear representation \(\varrho\) of G on the Banach space \(\mathcal{C}(\mathrm{G})\) defined by

\[
\varrho (g) f (x) = f (g ^ {- 1} x g)
\]

for all \(f \in \mathcal{C}(\mathrm{G})\) and every \(x \in \mathrm{G}\) . It is a continuous isometric representation and \(\mathcal{C}(\mathrm{G}^{\sharp})\) is the space of invariant elements of this representation. The continuous projector \(p = \varrho(1)\) of \(\mathcal{C}(\mathrm{G})\) therefore has as its image \(\mathcal{C}(\mathrm{G}^{\sharp})\) (cor. 3 of V, p. 466); it has norm \(\leqslant 1\) .

Let \(f \in \mathcal{C}(\mathrm{G}^{\sharp})\) and let \(\varepsilon > 0\) . There exists \(\widetilde{f} \in \Theta(\mathrm{G})\) such that \(\|f - \widetilde{f}\| \leqslant \varepsilon\) (prop. 4 of V, p. 462), and we then have \(\|f - p(\widetilde{f})\| = \|p(f - \widetilde{f})\| \leqslant \varepsilon\) ; since \(p(\widetilde{f}) \in \mathcal{C}(\mathrm{G}^{\sharp})\) , we conclude that \(\Theta(\mathrm{G}^{\sharp})\) is dense in \(\mathcal{C}(\mathrm{G}^{\sharp})\) .

We denote by R(G) the Grothendieck ring of continuous representations of G in separated finite-dimensional complex vector spaces (cf. LIE, IX, p. 70 and A, VIII, p. 182, applied to the additive class of continuous representations of G in separated finite-dimensional complex vector spaces, viewed as \(\mathbf{C}[G]\)-modules). Since every continuous linear representation of G in a finite-dimensional separated topological vector space is semisimple (prop. 1 of V, p. 457), the abelian group R(G) is free, and the classes of irreducible unitary representations \(\pi \in \widehat{\mathbf{G}}\) form a basis (A, VIII, p. 186, prop. 7).

COROLLARY 2. --- a) The family of characters of the irreducible representations of G is a basis of \(\Theta(G^{\sharp})\) ;

\begin{enumerate}
\def\labelenumi{\alph{enumi})}
\setcounter{enumi}{1}
\tightlist
\item
  The map \(u: \mathrm{R}(\mathrm{G}) \otimes_{\mathbf{Z}} \mathbf{C} \to \Theta(\mathrm{G}^{\sharp})\) such that \(u(\pi \otimes 1) = \chi_{\pi}\) for every \(\pi \in \widehat{\mathrm{G}}\) is an isomorphism of algebras over \(\mathbf{C}\) .
\end{enumerate}

The first assertion follows from Proposition 9 and Corollary 1.

Since the classes of the representations \(\pi \in \widehat{\mathbf{G}}\) form a basis of the free \(\mathbf{Z}\)-module R(G), the map \(u\) is well defined. It is a morphism of \(\mathbf{C}\) -algebras, and it is an isomorphism by \(a\) ).

COROLLARY 3. --- Let H be a locally compact topological group such that G is a compact subgroup of H. Let \(\varrho\) be a unitary representation of H in a Hilbert space E. If the \(\pi\)-isotypic component of the restriction of \(\varrho\) to G is finite-dimensional for all \(\pi \in \widehat{\mathbf{G}}\) , then the representation \(\varrho\) of H is semisimple and every irreducible unitary representation of H has finite multiplicity in \(\varrho\) .

Denote by \(\sigma\) the restriction of \(\varrho\) to the compact subgroup G of H. Let \(\pi\) be an irreducible representation of G. The endomorphism \(\sigma(\chi_{\pi}) \in \mathcal{L}(\mathrm{E})\) has finite rank, since its image is the \(\overline{\pi}\)-isotypic component of \(\sigma\) (cor. 2 of V, p. 465) and it is of finite dimension by hypothesis. Consequently, the endomorphism \(\sigma(f)\) has finite rank for every \(f \in \Theta(\mathrm{G}^{\sharp})\) and is compact for every \(f \in \mathcal{C}(\mathrm{G}^{\sharp})\) (cor. 1 of V, p. 468 and cor. of prop. 2 of III, p. 4).

Denote by \(j\) the canonical injection from G into H. This is a continuous map that is \(\nu\) -proper for every measure \(\nu\) on G (INT, V, p. 69, § 6, no. 1, remark 1). Let \(\mathfrak{B}\) be the filter of compact neighborhoods of the element \(e\) in G. For every V \(\in\) \(\mathfrak{B}\) there exists a continuous positive central function \(f_{\mathrm{V}}\) on G with support contained in V whose integral over G equals 1. Let \(\beta_{\mathrm{V}}\) be the image measure \(j(f_{\mathrm{V}} \cdot \mu)\) ; it is a positive measure with compact support on H such that \(\varrho(\beta_{\mathrm{V}}) = \sigma(f_{\mathrm{V}})\) (INT, V, p. 71, § 6, no. 2, th. 1). According to what precedes, the filter base on \(\mathcal{M}_{c}(\mathrm{H})\) obtained as the image of \(\mathfrak{B}\) under the map \(\mathrm{V} \mapsto \beta_{\mathrm{V}}\) satisfies the conditions of Proposition 2 of V, p. 402, and the assertion follows.

COROLLARY 4 (Weyl's equidistribution criterion)

Let M be a set of central positive measures on G such that \(\nu(G)=1\) for all \(\nu\in M\) . Let \(\mathfrak{F}\) be a filter on M. For the filter \(\mathfrak{F}\) to converge vaguely to the measure \(\mu^{\sharp}\) on \(G^{\sharp}\) , it is necessary and sufficient that, for every nontrivial irreducible unitary representation \(\pi\) , one has

\[
\lim _ {\nu , \mathfrak {F}} \int_ {\mathrm{G}} \chi_ {\pi} (x) d \nu (x) = 0.
\]

Since \(\nu (\mathrm{G}^{\sharp}) = 1 = \mu^{\sharp}(\mathrm{G}^{\sharp})\) for \(\nu \in \mathrm{M}\) , the hypothesis means that

\[
\lim _ {\nu , \mathfrak {F}} \int_ {\mathrm{G} ^ {\sharp}} \chi_ {\pi} (x) d \nu (x) = \int_ {\mathrm{G} ^ {\sharp}} \chi_ {\pi} (x) d \mu^ {\sharp} (x)
\]

for every representation \(\pi\in\widehat{\mathrm{G}}\) , hence it is equivalent to the condition

\[
\lim _ {\nu , \mathfrak {F}} \int_ {\mathrm{G} ^ {\sharp}} f (x) d \nu (x) = \int_ {\mathrm{G} ^ {\sharp}} f (x) d \mu^ {\sharp} (x)
\]

for every function \(f \in \Theta(\mathrm{G}^{\sharp})\) by linearity. Since the space \(\Theta(\mathrm{G}^{\sharp})\) is dense in \(\mathcal{C}(\mathrm{G}^{\sharp})\) (corollary 1), the hypothesis is therefore equivalent to the convergence of the filter \(\mathfrak{F}\) to \(\mu^{\sharp}\) in \(\mathcal{M}(\mathrm{G}^{\sharp})\) endowed with the topology of pointwise convergence in \(\mathcal{C}(\mathrm{G})\) , which coincides with the topology of vague convergence since G is compact (cf. INT, III, p. 59, § 1, no. 9).

